% GENERATED FILE. Edit canonical lessons, then run npm run generate:books.
% canonical-source-hash: fnv1a64:2f533c3f8346daae
% canonical-lessons: KA-C98-boredom, KA-C98-enthusiasm, KA-C98-shyness, KA-C98-alarm, KA-C98-peace

\chapter{Moods}
\label{ch:ka-moods}
\hlchaptermodality{98}

\begin{chapteropening}
\textbf{By the end of this chapter:} \emph{I can name boredom, enthusiasm, shyness, alarm and peace of mind.}

Complete the last lesson of chapter 98: i can name boredom, enthusiasm, shyness, alarm and peace of mind.
\end{chapteropening}

\section[bēsara]{\kn{ಬೇಸರ} (bēsara) — boredom, weariness of spirit}

\label{lesson:KA-C98-boredom}



\begin{quote}
Before the new one: say the Kannada for don't forget, then the Kannada for kindness.
\end{quote}

\subsection*{You'll want to know: \kn{ಬೇಸರ}}
\textbf{\kn{ಬೇಸರ}} (\emph{bēsara}) — \textquotedblleft{}boredom\textquotedblright{}, and the low mood that comes with it.

\textbf{\kn{ನನಗೆ} \kn{ಬೇಸರ} \kn{ಆಗಿದೆ}} (\emph{nanage bēsara āgide}) — \textquotedblleft{}I'm bored\textquotedblright{}, with the person in \textbf{\kn{ನನಗೆ}} once more.

\begin{grammarlens}[title={What you've built}]
A mood, said like hunger and tiredness.
\end{grammarlens}

\subsection*{Your turn}
\begin{itemize}
\raggedright
  \item \emph{Say it:} \emph{bēsara}
  \item \emph{Say it:} \emph{bēsara}, once more
  \item \emph{Say it:} say \emph{nanage bēsara āgide}
  \item \emph{Recall:} say the Kannada for don't forget, then the Kannada for kindness, then say \emph{bēsara} again
\end{itemize}

\begin{tcolorbox}[breakable,colback=teal!4,colframe=teal!35!black,title={Before you move on}]
What does \kn{ಬೇಸರ} mean? (\textquotedblleft{}Boredom, weariness of spirit\textquotedblright{}.) Say it once more.
\end{tcolorbox}
\section[utsāha]{\kn{ಉತ್ಸಾಹ} (utsāha) — enthusiasm, eagerness}

\label{lesson:KA-C98-enthusiasm}



\begin{quote}
Before the new one: say the Kannada for kindness, then the Kannada for boredom.
\end{quote}

\subsection*{You'll want to know: \kn{ಉತ್ಸಾಹ}}
\textbf{\kn{ಉತ್ಸಾಹ}} (\emph{utsāha}) — \textquotedblleft{}enthusiasm, eagerness\textquotedblright{}, the opposite of \textbf{\kn{ಬೇಸರ}}.

\begin{grammarlens}[title={What you've built}]
The mood that gets things started.
\end{grammarlens}

\subsection*{Your turn}
\begin{itemize}
\raggedright
  \item \emph{Say it:} \emph{utsāha}
  \item \emph{Say it:} \emph{utsāha}, once more
  \item \emph{Say it:} say \emph{utsāha}
  \item \emph{Recall:} say the Kannada for kindness, then the Kannada for boredom, then say \emph{utsāha} again
\end{itemize}

\begin{tcolorbox}[breakable,colback=teal!4,colframe=teal!35!black,title={Before you move on}]
What does \kn{ಉತ್ಸಾಹ} mean? (\textquotedblleft{}Enthusiasm, eagerness\textquotedblright{}.) Say it once more.
\end{tcolorbox}
\section[nācike]{\kn{ನಾಚಿಕೆ} (nācike) — shyness, embarrassment}

\label{lesson:KA-C98-shyness}



\begin{quote}
Before the new one: say the Kannada for boredom, then the Kannada for enthusiasm.
\end{quote}

\subsection*{You'll want to know: \kn{ನಾಚಿಕೆ}}
\textbf{\kn{ನಾಚಿಕೆ}} (\emph{nācike}) — \textquotedblleft{}shyness, embarrassment\textquotedblright{}.

\textbf{\kn{ನಾಚಿಕೆ} \kn{ಬೇಡ}} (\emph{nācike bēḍa}) — \textquotedblleft{}don't be shy\textquotedblright{}, a kind thing to say to a new learner.

\begin{grammarlens}[title={What you've built}]
A feeling every learner knows.
\end{grammarlens}

\subsection*{Your turn}
\begin{itemize}
\raggedright
  \item \emph{Say it:} \emph{nācike}
  \item \emph{Say it:} \emph{nācike}, once more
  \item \emph{Say it:} say \emph{nācike bēḍa}
  \item \emph{Recall:} say the Kannada for boredom, then the Kannada for enthusiasm, then say \emph{nācike} again
\end{itemize}

\begin{tcolorbox}[breakable,colback=teal!4,colframe=teal!35!black,title={Before you move on}]
What does \kn{ನಾಚಿಕೆ} mean? (\textquotedblleft{}Shyness, embarrassment\textquotedblright{}.) Say it once more.
\end{tcolorbox}
\section[gābari]{\kn{ಗಾಬರಿ} (gābari) — alarm, panic}

\label{lesson:KA-C98-alarm}



\begin{quote}
Before the new one: say the Kannada for enthusiasm, then the Kannada for shyness.
\end{quote}

\subsection*{You'll want to know: \kn{ಗಾಬರಿ}}
\textbf{\kn{ಗಾಬರಿ}} (\emph{gābari}) — \textquotedblleft{}alarm, panic, fright\textquotedblright{}.

\textbf{\kn{ಗಾಬರಿ} \kn{ಆಗಬೇಡಿ}} (\emph{gābari āgabēḍi}) — \textquotedblleft{}don't panic\textquotedblright{}.

\begin{grammarlens}[title={What you've built}]
A sudden fright, and how to calm it.
\end{grammarlens}

\subsection*{Your turn}
\begin{itemize}
\raggedright
  \item \emph{Say it:} \emph{gābari}
  \item \emph{Say it:} \emph{gābari}, once more
  \item \emph{Say it:} say \emph{gābari āgabēḍi}
  \item \emph{Recall:} say the Kannada for enthusiasm, then the Kannada for shyness, then say \emph{gābari} again
\end{itemize}

\begin{tcolorbox}[breakable,colback=teal!4,colframe=teal!35!black,title={Before you move on}]
What does \kn{ಗಾಬರಿ} mean? (\textquotedblleft{}Alarm, panic\textquotedblright{}.) Say it once more.
\end{tcolorbox}
\section[nemmadi]{\kn{ನೆಮ್ಮದಿ} (nemmadi) — peace of mind}

\label{lesson:KA-C98-peace}



\begin{quote}
Before the new one: say the Kannada for shyness, then the Kannada for alarm.
\end{quote}

\subsection*{You'll want to know: \kn{ನೆಮ್ಮದಿ}}
\textbf{\kn{ನೆಮ್ಮದಿ}} (\emph{nemmadi}) — \textquotedblleft{}peace of mind, contentment\textquotedblright{}, the calm after worry.

\begin{grammarlens}[title={What you've built}]
The quiet end of the chapter.
\end{grammarlens}

\subsection*{Your turn}
\begin{itemize}
\raggedright
  \item \emph{Say it:} \emph{nemmadi}
  \item \emph{Say it:} \emph{nemmadi}, once more
  \item \emph{Say it:} say \emph{nemmadi}
  \item \emph{Recall:} say the Kannada for shyness, then the Kannada for alarm, then say \emph{nemmadi} again
\end{itemize}

\begin{tcolorbox}[breakable,colback=teal!4,colframe=teal!35!black,title={Before you move on}]
What does \kn{ನೆಮ್ಮದಿ} mean? (\textquotedblleft{}Peace of mind\textquotedblright{}.) Say it once more.
\end{tcolorbox}
